\documentclass[11pt,a4paper]{article}
\usepackage[margin=25mm]{geometry}
\usepackage{amsmath,amssymb,amsthm,hyperref}
\newtheorem{theorem}{Theorem}
\newtheorem{lemma}[theorem]{Lemma}
\newtheorem{corollary}[theorem]{Corollary}
\newcommand{\F}{\mathbb F}
\newcommand{\dist}{\operatorname{dist}}
\title{Graphs of girth five and diameter four}
\author{Human Compression project}
\date{}
\begin{document}\maketitle

We give complete proofs of the general results used in the recovery of
Araujo-Pardo, Conder, Garc\'ia-Col\'in, Kiss and Leemans,
\emph{A note on girth-diameter cages}~\cite{source}.
The lower bound and the middle-graph construction are their results.
The use of biaffine planes and the addition of edges within their deleted
point and line classes are earlier constructions, attributed in that
paper to Brown and subsequent authors. Our contribution here is an
explicit common proof with precise endpoint conditions and a criterion
that verifies the diameter, in addition to the girth. We make no priority
claim for the constructions.

Graphs are finite, simple and undirected. A graph of diameter four is
connected. For a graph which need not be connected, distance between
different components is infinity; a graph without cycles has infinite
girth. The source's finite classifications require the separate complete
census and are not consequences of these written theorems alone.

\section{The middle graph}

\begin{theorem}\label{thm:middle}
Let $k\geq3$. A $k$-regular graph of girth at least five and diameter
at least four has at least $k^2+k+2$ vertices. Equality is attained
precisely by the following construction.

Put $V_k=\{(i,j):1\leq i,j\leq k,\ i\ne j\}$. Choose a simple
$(k-2)$-regular graph $H$ on $V_k$ such that:
\begin{enumerate}
\item $H$ has no triangle or quadrangle;
\item any two different vertices in the same row or the same column
are neither adjacent nor joined by a path of length two in $H$.
\end{enumerate}
Adjoin vertices $r,c,r_1,\ldots,r_k,c_1,\ldots,c_k$ and edges
\[
 rr_i,\qquad r_i(i,j),\qquad cc_j,\qquad c_j(i,j)
 \quad(i\ne j).
\]
The resulting graph has exactly girth five and diameter four.
Conversely, every graph attaining the bound is obtained in this way
from any chosen ordered pair of vertices at distance at least four.
\end{theorem}

\begin{proof}
Choose $r,c$ at distance at least four. Their neighbourhoods $R,C$
are disjoint independent sets of size $k$, and there is no edge
between $R$ and $C$. The $k(k-1)$ edges from $R$ other than those
to $r$ have distinct other endpoints: a repeated endpoint would
give a quadrangle through $r$. These endpoints lie outside
$\{r,c\}\cup R\cup C$. Consequently the set $M$ of remaining
vertices has size at least $k(k-1)$, which proves the bound.

Suppose equality holds. Every vertex of $M$ has exactly one neighbour
in $R$; the same argument from $c$ gives exactly one in $C$.
Write $R=\{r_i\}$, $C=\{c_j\}$ and let $R_i,C_j$ be their
respective sets of neighbours in $M$. Each has size $k-1$.
The sets $R_i$ partition $M$, as do the sets $C_j$, and
$|R_i\cap C_j|\leq1$: two common vertices would form a quadrangle
with $r_i,c_j$. Thus the $k$ by $k$ matrix of these intersection
sizes is a zero-one matrix with every row and column sum $k-1$.
Its zero positions form a permutation matrix. Relabelling the
columns places the zeros on the diagonal and identifies $M$ with
$V_k$. Each middle vertex has precisely two neighbours outside
$M$, so the induced graph $H$ has degree $k-2$.

Triangles and quadrangles in $H$ would persist in the original
graph. An edge between two vertices of $R_i$, or a path of length
two between them, would form a triangle or quadrangle with $r_i$;
the column statement follows in the same way. This proves necessity
of the two conditions. Notice that the disjointness of two different
$R_i$ follows from the absence of quadrangles through $r$, rather
than from the absence of triangles in $H$.

Conversely, carry out the stated construction. A middle vertex has
degree $(k-2)+2=k$, each $r_i,c_j$ has degree $1+(k-1)=k$, and
$r,c$ also have degree $k$. There are $k(k-1)+2k+2=k^2+k+2$
vertices. Every middle vertex is at distance two from each of $r,c$;
any two middle vertices are therefore at distance at most four.
A middle vertex is within three of every $r_i$ or $c_j$, by going
through $r$ or $c$. Two row vertices, or two column vertices, are
within two. If $i\ne j$, then $r_i,(i,j),c_j$ is a path of length
two. For $i=j$, choose $h\ne i$ and use
$r_i,(i,h),c_h,c,c_i$, of length four. Finally $r$ is within three
of every column vertex, and $c$ within three of every row vertex.
There is a path of length four from $r$ to $c$ through any middle
vertex. There is no shorter one, since $R\cap C=\varnothing$ and
there are no edges between $R$ and $C$. The diameter is exactly four.

There can be no triangle or quadrangle through $r$: a triangle
would require an edge within $R$, while a quadrangle would require
a middle vertex adjacent to two row vertices. The same holds for
$c$. A short cycle avoiding both $r,c$ is either contained in $H$
or contains row or column vertices, which are pairwise nonadjacent.
With just one such vertex it violates the row or column distance
condition. With two it would be a quadrangle. Two row vertices
cannot have a common middle neighbour, nor can two column vertices;
one row and one column vertex have at most one common neighbour.
This excludes all remaining short cycles.

It remains to prove the existence of a pentagon, not merely the
absence of shorter cycles. Since $k\geq3$, the graph $H$ has an
edge $(i,j)(i',j')$. Its endpoints have $i\ne i'$. Then
\[
 r,r_i,(i,j),(i',j'),r_{i'},r
\]
is a cycle of length five. This completes the proof.
\end{proof}

The separate diameter condition on $H$ in the source is redundant:
a row contains $k-1\geq2$ vertices, and any two of them have distance
at least three, possibly infinity. This convention includes the
disconnected middle graph $3K_2$ when $k=3$.
The restriction $k\geq3$ is necessary. When $k=2$, the prescribed
middle graph has two isolated vertices and the extension is $C_8$,
whose girth is eight. In fact a connected 2-regular graph is a cycle,
so no graph can have degree two, girth five and diameter four.

\begin{lemma}\label{lem:labels}
Every graph attaining the bound in Theorem~\ref{thm:middle} occurs in
an enumeration of all labelled graphs $H$ satisfying its two conditions.
Once the ordered vertices $r,c$ have been chosen, the remaining
freedom in the grid labelling is a simultaneous permutation of the
row and column indices. Interchanging $r,c$ transposes the grid.
To classify graphs without distinguished vertices, one must in
addition identify outputs by isomorphisms of the entire extended graph.
\end{lemma}
\begin{proof}
The construction in the necessity proof applies to every chosen
ordered pair. Labelling its $k$ row neighbours leaves exactly one
missing column neighbour for each row, so the diagonal convention
then forces all column labels. This proves the description of the
remaining freedom. An isomorphism need not preserve the selected
pair $r,c$; two different choices in the same graph may give grids
not related by these labelling changes. Complete graph isomorphism
therefore supplies the final equivalence relation, and cannot be
replaced by the simultaneous-index action alone.
\end{proof}

\section{Biaffine incidence graphs}

\begin{lemma}\label{lem:planes}
In a projective plane of order $q\geq3$, delete a point $P$, a line
$\ell$, all points on $\ell$ and all lines through $P$. The remaining
incidence graph is $q$-regular, has girth six and diameter four.
Its order is $2q^2$ if $P\in\ell$ and $2(q^2-1)$ otherwise.
Two vertices have distance four precisely when they are different
points on one deleted line or different lines through one deleted point.
\end{lemma}
\begin{proof}
There are $q^2+q+1$ points and as many lines. If $P\in\ell$ the
deleted point and line sets each have $q+1$ members; otherwise each
has $q+2$ members. A retained point loses exactly its line to $P$,
and a retained line loses exactly its point on $\ell$, leaving
degree $q$ on either side.

Two retained points on different lines through $P$ have their unique
joining line still present, hence distance two. Two retained points
on the same such line have no retained common neighbour. Choose
a retained point on a different line through $P$; the two joining
lines give a path of length four. The dual statement holds for lines.
For a nonincident retained point $A$ and line $b$, choose one of the
$q$ retained points $B$ on $b$ other than its possible intersection
with $PA$. Then $AB$ is retained, giving a path of length three.
Thus these are all the distances. A deleted line has at least two
retained points when $q\geq3$, so the diameter is four.

The graph is bipartite and has no quadrangle because joining lines
are unique. To find a hexagon, choose retained points $A,B$ on
different lines through $P$, and a third line through $P$ distinct
from these and from $\ell$ if $P\in\ell$. This line contains at
least two retained points, at most one on $AB$. Choose another,
say $C$. The three joining lines of the noncollinear points $A,B,C$
are retained and yield a hexagon.
\end{proof}

For $q=2$ the two incidence graphs are respectively $C_8$ and $C_6$.
Indeed the same distance argument gives connected 2-regular graphs
on eight and six vertices. Thus the source's preliminary girth and
diameter statement needs these small-order exceptions.

We now use coordinates to verify two standard ways of adding edges.
These proofs retain only elementary finite-field facts, including
cyclicity of the multiplicative group.

\begin{theorem}\label{thm:additive}
Let $q$ be a prime power, and let $S,T\subseteq\F_q\setminus\{0\}$
be nonempty sets closed under negation. On vertices
$P(x,y),L(m,b)$, with $x,y,m,b\in\F_q$, join
\[
 P(x,y)\sim L(m,b)\quad\Longleftrightarrow\quad y=mx+b.
\]
Also join $P(x,y)$ to $P(x,y+s)$ for $s\in S$, and $L(m,b)$ to
$L(m,b+t)$ for $t\in T$. Suppose $S\cap T=\varnothing$ and both
additive Cayley graphs defined by $S,T$ have no triangles or quadrangles.
The resulting graph has girth five and diameter at most four.
It is $(q+a)$-regular when $|S|=|T|=a$ and has order $2q^2$.

Write $2S=S+S$ and $3S=S+S+S$, and similarly for $T$.
Its diameter is four if and only if at least one of
\[
 \F_q\setminus(\{0\}\cup S\cup2S\cup3S\cup T),\qquad
 \F_q\setminus(\{0\}\cup T\cup2T\cup3T\cup S)
\]
is nonempty.
\end{theorem}

\begin{proof}
The original incidence graph is that of the first biaffine plane:
take $\ell$ to be the line at infinity and $P$ the vertical point
on it. Its two parts have $q^2$ vertices each. Two point vertices
with different first coordinates have exactly one common line
vertex. Two different point vertices with the same first coordinate
have none and have distance four. The dual statement holds for lines
of equal or unequal slopes. A nonincident point and line have
distance three: choose a line through the point of slope different
from the target line and take their intersection. These facts also
hold at $q=2$ and prove the upper diameter bound after edges are added.

An added point edge joins two vertices in the same vertical class.
They cannot have a common line neighbour, so it creates no triangle
with incidence edges. The same holds dually. An entirely internal
triangle is excluded by the Cayley hypothesis.
In a quadrangle the number of incidence edges is even. Four such
edges are impossible in the original incidence graph. Zero are
excluded within the fibres. If two internal edges lie on the same
side, their endpoints would be distinct vertices of one fibre with
a common incidence neighbour, also impossible. The remaining case
has one internal point edge and one internal line edge, necessarily
with a common $x$ and a common $m$. Subtracting the two incidence
equations gives $s=t$ for some $s\in S,t\in T$, contrary to
disjointness. Thus the girth is at least five. The endpoints of any
added edge had distance four in the incidence graph; adjoining it
to a shortest such path produces a pentagon.

Only pairs within a point fibre or a line fibre could still be at
distance four; all other distances were at most three already.
For $P(x,y)$ and $P(x,y+d)$, a path of length at most three with
no incidence edges changes the second coordinate by an element of
$\{0\}\cup S\cup2S\cup3S$. A path with two incidence edges and
one internal edge changes it by an element of $S$ if that edge is
on the point side. If it is on the line side, the change is in $T$,
by subtracting the incidence equations. Two incidence edges alone
can return to the same point fibre only at the original vertex.
These exhaust the possibilities by parity of the side changes.
Every listed change is conversely realized by a walk of length
at most three; deleting repetitions gives a path. The point-fibre
distance criterion follows. Interchanging the two sides gives the
other displayed set and proves the assertion.
\end{proof}

\begin{corollary}\label{cor:additive}
For $q=2^r\geq4$ there is a $(q+1)$-regular graph of girth five,
diameter four and order $2q^2$. For $q=p^r>7$ with prime $p\geq5$,
there is a $(q+2)$-regular graph with these same girth, diameter and order.
\end{corollary}
\begin{proof}
In characteristic two choose different nonzero $a,b$, and put
$S=\{a\}$, $T=\{b\}$. The internal graphs are matchings, and
$a+b$ is outside $\{0,a,b\}$, which is the entire set of changes
of length at most three in the point fibre. The theorem applies.

In odd characteristic put $S=\{1,-1\}$ and $T=\{2,-2\}$.
They are disjoint for $p\geq5$, and each internal graph is a disjoint
union of cycles of length $p$, hence has no triangle or quadrangle.
The possible short changes lie among $0,\pm1,\pm2,\pm3$.
There are at most seven of them, fewer than $q$, so the diameter is four.
\end{proof}

The first construction gives the source's degree-five, order-32
example at $q=4$ and its quoted extension in characteristic two.
The second also explains why a cycle length of seven does not
alone ensure diameter four: for $q=7$ these short changes fill the
field and the same construction has diameter at most three.
A 7-cycle itself has diameter three. We therefore do not carry over
the source's illustrative ``$q\geq7$'' diameter assertion.

\begin{theorem}\label{thm:multiplicative}
For every prime power $q\geq11$ there is a $(q+2)$-regular graph
of girth five, diameter four and order $2(q^2-1)$.
\end{theorem}
\begin{proof}
Use disjoint copies of $\F_q^2\setminus\{0\}$ for point vectors
$v$ and line covectors $u$, with incidence $u(v)=1$. This is the
second biaffine plane: delete the origin and the line at infinity.
An independent pair of point vectors has exactly one common
covector satisfying these equations; a distinct proportional pair
has none. Choosing an independent third vector joins a proportional
pair by a path of length four. Dually the same holds for covectors.
For nonincident $v,u$, the affine line $u(w)=1$ has $q$ vectors,
at most one proportional to $v$. Choose another $w$ and take the
unique common covector of $v,w$, giving a path of length three.

Let $\lambda$ generate $\F_q^\times$, of order $N=q-1$. Add
edges $v\sim\lambda v$ and $u\sim\lambda^2u$, with their inverses.
The internal point graphs are cycles of length $N$, and the line
graphs cycles of length $N/\gcd(N,2)$. Both lengths are at least
five. Each side acquires degree two, so the total degree is $q+2$.

There is no triangle involving two incidence edges: two distinct
proportional vectors cannot have a common incidence neighbour.
Internal triangles are excluded by the cycle lengths. As in the
proof of Theorem~\ref{thm:additive}, a quadrangle with two internal
edges on one side would give the same impossible common neighbour.
With one internal edge on each side, write their multipliers as
$a\in\{\lambda,\lambda^{-1}\}$ and
$b\in\{\lambda^2,\lambda^{-2}\}$. The two incidence conditions
would give $ab=1$, impossible for $N\geq10$.
The original incidence graph has no quadrangle, and the internal
cycles have none. Thus the girth is at least five. Any added edge,
together with an original path of length four between its endpoints,
gives a pentagon.

The diameter is at most four by the original incidence paths.
Consider $v$ and $\lambda^4v$. A path of length at most three
staying on the point side can multiply $v$ only by $\lambda^j$
with $-3\leq j\leq3$. A path with two incidence edges and one
internal point edge gives $\lambda$ or $\lambda^{-1}$ as multiplier.
If its internal edge is on the line side, the two equations force
the point multiplier to be $\lambda^2$ or $\lambda^{-2}$.
Two incidence edges alone cannot join distinct proportional vectors.
Since $N\geq10$, none of these multipliers is $\lambda^4$.
The selected pair therefore has distance four.
\end{proof}

In particular Theorem~\ref{thm:multiplicative} supplies the quoted
type-two family for all powers of two $q\geq16$. Its explicit
parameter range is part of the claim. In the cycle construction
at $q=8$, every class of proportional vectors is a 7-cycle, so its
vertices are within distance three; all other pairs were already
within three. That construction consequently cannot justify
diameter four at $q=8$.

The two families imply the principal asymptotic statement of the
source: $n_0(k;5,4)\leq2k^2$ for infinitely many $k$. Indeed
Corollary~\ref{cor:additive} already provides this with
$k=2^r+1$, $r\geq2$. They do not settle the problem of attaining
the lower bound $k^2+k+2$ for further degrees.

\begin{thebibliography}{9}
\bibitem{source}
G. Araujo-Pardo, M. Conder, N. Garc\'ia-Col\'in, Gy. Kiss and D. Leemans,
\emph{A note on girth-diameter cages}, arXiv:2401.15539v1,
28 January 2024, manuscript dated 30 January 2024, Sections 1--5.
The project registry uses the later publication year 2025.
\end{thebibliography}
\end{document}
